\scanpage{36}
\mspage{17}
est connu quand on connaît
\[
(42)\qquad\Pi_1(f):\Pi_1(G')\longrightarrow\Pi_1(G^{\deleted{\unsure{'}}}),
\]
puisqu'alors on connaît l'action de $f$ sur $S$ via l'action de $\Pi_1(f)$ sur les objets, et l'action de $f$ sur les arcs via l'action de $\Pi_1(f)$ sur les flèches, grâce au diagramme comm.

(43)\quad\diagram{\begin{tikzcd}
\widetilde A(G')\arrow[r,"\widetilde A(f)"]\arrow[d,hook]&\widetilde A(G)\arrow[d,hook]\\
\operatorname{Fl}(\Pi_1(G'))\arrow[r,"\operatorname{Fl}(\Pi_1(f))"]&\operatorname{Fl}(\Pi_1(G))
\end{tikzcd}}.

En fait, les morphismes de $G$ dans $G'$ s'identifient aux morphismes de groupoïdes $\Pi_1(G')$ dans $\Pi_1(G)$, qui transforment arcs \inserted{de $G'$} en arcs de $G$.

Ceci étant, on généralise la notion de morphisme de graphes en celle de \uline{morphisme généralisé} de $G'$ dans $G^{\deleted{\unsure{'}}}\deleted{\unsure{,}}$, (ou ``\uline{morphisme de générisation}'' de $G'$ dans $G^{\deleted{\unsure{'}}}$)
\mspage{35}
